\documentclass[12pt,a4paper]{article}

% ---------------------------------------------------------------------------
% Formatação conforme a "Alternativa" explicitamente permitida pelo enunciado
% (Seção 6): A4, coluna simples, fonte Times 12pt, margens 3,5cm (superior),
% 2,5cm (inferior), 3,0cm (laterais), sem cabeçalho/rodapé/numeração de página.
% Caso o grupo prefira o template oficial da SBC via Overleaf
% (https://www.overleaf.com/latex/templates/sbc-conferences-template/blbxwjwzdngr),
% basta substituir o preâmbulo abaixo pela classe sbc.cls baixada de lá — todo o
% conteúdo das seções foi escrito para ser independente da classe usada.
% ---------------------------------------------------------------------------

\usepackage[a4paper,top=3.5cm,bottom=2.5cm,left=3.0cm,right=3.0cm]{geometry}
\usepackage[utf8]{inputenc}
\usepackage[T1]{fontenc}
\usepackage[brazil,english]{babel}
\usepackage{mathptmx}     % fonte Times
\usepackage{graphicx}
\usepackage{booktabs}
\usepackage{amsmath}
\usepackage{hyperref}
\usepackage{url}
\usepackage{caption}
\usepackage{titlesec}

% Compactação tipográfica leve (não altera fonte/margens exigidas pelo enunciado):
% reduz espaçamento antes/depois de títulos de seção e de legendas de tabela/figura,
% para aproveitar melhor o limite rígido de 4 páginas.
\titlespacing*{\section}{0pt}{8pt}{4pt}
\titlespacing*{\subsection}{0pt}{6pt}{2pt}
\captionsetup{skip=4pt}
\setlength{\textfloatsep}{8pt plus 2pt minus 2pt}
\setlength{\intextsep}{8pt plus 2pt minus 2pt}

\pagestyle{empty}         % sem cabeçalho, rodapé ou numeração de página

\title{\textbf{Baseline Clássico para Detecção de Embolia Pulmonar em Angio-TC de Tórax: Engenharia de Características e Modelos Não Profundos sobre o Desafio RSNA 2020}}

\author{
Rawenne da Silva Leite, Wheverson de Abreu Lima, João Victor do Nascimento Silva
}

\date{}

\begin{document}
\selectlanguage{brazil}
\maketitle
\thispagestyle{empty}

\begin{center}
\small
\textit{Tópicos Especiais em Sistemas de Informação --- Curso de Sistemas de Informação}\\
\textit{Grupo G2 --- Desafio RSNA 2020: Pulmonary Embolism Detection}
\end{center}

\vspace{0.3cm}

\noindent\textbf{Resumo.} Este trabalho constrói e avalia uma solução \emph{baseline}
para a detecção de embolia pulmonar (EP) em exames de angiotomografia computadorizada
de tórax, restrita exclusivamente a métodos clássicos de visão computacional:
engenharia de características manuais (textura, forma, gradiente, intensidade e
radiômica) combinada a classificadores rasos (regressão logística, SVM, $k$-NN,
Random Forest, \emph{gradient boosting} e MLP raso). O objetivo não é maximizar
desempenho, mas estabelecer o piso metodologicamente rigoroso contra o qual
abordagens futuras baseadas em aprendizado profundo deverão ser comparadas.
Descrevemos o protocolo experimental --- particionamento por paciente, validação
cruzada estratificada e agrupada, métricas apropriadas a um problema clinicamente
desbalanceado --- e discutimos as dificuldades intrínsecas do desafio: rótulos
hierárquicos, esparsidade espacial do achado ao longo do volume e variabilidade de
aquisição entre instituições. O repositório de código, incluindo um pipeline
validado de ponta a ponta, está disponível publicamente (link na nota de rodapé).

\vspace{0.2cm}

\selectlanguage{english}
\noindent\textbf{Abstract.} We build and evaluate a baseline solution for pulmonary
embolism (PE) detection in chest CT pulmonary angiography, restricted exclusively to
classical computer vision methods: hand-crafted feature engineering (texture, shape,
gradient, intensity and radiomics) combined with shallow classifiers (logistic
regression, SVM, $k$-NN, Random Forest, gradient boosting and a shallow MLP). The goal
is not to maximize performance but to establish a methodologically rigorous floor
against which future deep-learning approaches must be compared. We describe the
experimental protocol --- patient-level partitioning, stratified grouped
cross-validation, metrics appropriate to a clinically imbalanced problem --- and
discuss the challenges intrinsic to the task: hierarchical labels, spatial sparsity of
the finding along the volume, and cross-institution acquisition variability. The code
repository, including an end-to-end validated pipeline, is publicly available (link
in the footnote).
\selectlanguage{brazil}

\footnotetext[0]{Repositório: \url{https://github.com/silvarawenne/tp1_g2} (branch \texttt{claude/laughing-rubin-79l5wt}). \textbf{Uso de IA generativa:} partes do texto e do código foram redigidas com apoio de IA generativa (Claude); revisadas pelo grupo, que assume responsabilidade integral por sua correção, inclusive a verificação bibliográfica das referências.}

\section{Introdução}

A embolia pulmonar (EP) é uma condição potencialmente fatal, diagnosticada
primariamente por angiotomografia computadorizada de tórax (angio-TC). A leitura
manual de centenas de cortes por exame é demorada e sujeita a variabilidade
inter-observador, o que motiva sistemas de auxílio ao diagnóstico (CAD). O desafio
\emph{RSNA STR Pulmonary Embolism Detection} (2020) disponibilizou a maior base
pública anotada para esta tarefa até então \cite{colak2021rsnadataset}, com rótulos
hierárquicos no nível de corte e de exame.

A maioria dos trabalhos recentes recorre a redes neurais profundas
\cite{huang2020penet}. Este TP parte de premissa diferente, imposta pelo enunciado da
disciplina: construir um \emph{baseline} exclusivamente com métodos clássicos ---
características manuais combinadas a classificadores rasos --- estabelecendo, com
rigor metodológico, o piso de desempenho que uma futura abordagem profunda precisará
superar para justificar seu custo computacional adicional. Contribuições: (i) um
pipeline reprodutível de leitura DICOM, pré-processamento e extração de cinco
famílias de descritores clássicos; (ii) um protocolo sem vazamento de dados
(\texttt{StratifiedGroupKFold} por paciente, ajuste de hiperparâmetros em validação
aninhada); (iii) uma comparação sistemática descritor~$\times$~modelo contra o
baseline trivial obrigatório; (iv) uma discussão fundamentada das dificuldades
intrínsecas do desafio RSNA 2020.

\section{Trabalhos Relacionados}

A literatura de CAD para EP pré-datando o aprendizado profundo já explorava
características extraídas ao longo de estruturas vasculares segmentadas, com
classificadores de aprendizado multi-instância \cite{bi2007pemultipleinstance} --- a
mesma filosofia (características manuais + classificador clássico) adotada aqui,
ainda que sem a segmentação vascular supervisionada daquele trabalho. As famílias de
descritores usadas têm fundamentação estabelecida: textura via GLCM/Haralick
\cite{haralick1973textural} e LBP \cite{ojala2002lbp}; gradiente via HOG
\cite{dalal2005hog}; radiômica via PyRadiomics/IBSI \cite{vangriethuysen2017pyradiomics}.
Na modelagem, usamos os algoritmos clássicos mais estabelecidos para dados tabulares
desbalanceados: SVM \cite{cortes1995svm}, Random Forest
\cite{breiman2001randomforests}, \emph{gradient boosting} \cite{chen2016xgboost} e
SMOTE \cite{chawla2002smote}, frente ao desbalanceamento reportado na descrição
oficial do desafio \cite{colak2021rsnadataset}. Em contraste, trabalhos de ponta
neste desafio usam redes convolucionais 3D ponta a ponta, como o PENet
\cite{huang2020penet} --- este TP posiciona-se deliberadamente antes dessa
fronteira, fornecendo o número de referência contra o qual tais modelos deverão ser
comparados na fase seguinte do projeto.

\section{Metodologia}

\subsection{Base, amostragem e pré-processamento}
Utilizamos o conjunto \emph{RSNA STR Pulmonary Embolism Detection} (2020), hospedado
no Kaggle. Como os rótulos do teste oficial não são públicos, todo o protocolo é
construído sobre o conjunto de treino, com partição própria (\S\ref{sec:protocolo}).
Dado o volume da base (centenas de GB), trabalhamos com uma amostra estratificada por
classe e por paciente, com semente fixa e critério de inclusão documentados no
\texttt{README.md} -- a amostragem é parte do protocolo, não uma simplificação não
documentada. Cada corte DICOM é lido com \texttt{pydicom} e convertido para
Hounsfield Units (HU) via $\text{HU} = \text{pixel} \times \text{RescaleSlope} +
\text{RescaleIntercept}$. Aplicamos janelamento mediastinal/EP (centro 100~HU,
largura 700~HU), que realça o contraste entre o lúmen vascular opacificado e um
eventual defeito de enchimento, seguido de reamostragem isotrópica (1~mm/pixel) para
tornar os descritores comparáveis entre exames com \texttt{PixelSpacing} distintos,
recorte de região de interesse e redimensionamento final. Esta etapa é determinística
e por isso precede a divisão treino/teste sem risco de vazamento; normalizadores,
seletores, PCA e SMOTE, por outro lado, são encapsulados em
\texttt{sklearn.Pipeline} e ajustados somente na partição de treino de cada dobra.

\subsection{Extração de características e agregação corte~$\to$~exame}
Extraímos cinco famílias de descritores por corte: (i) \textbf{textura} --- GLCM/Haralick
\cite{haralick1973textural} e LBP uniforme \cite{ojala2002lbp}; (ii) \textbf{forma e
contorno} --- momentos de Hu e propriedades de região; (iii) \textbf{gradiente} ---
HOG \cite{dalal2005hog}; (iv) \textbf{intensidade} --- histograma e estatísticas de
ordem superior; (v) \textbf{radiômica} --- subconjunto IBSI via PyRadiomics
\cite{vangriethuysen2017pyradiomics}. Como a tarefa é volumétrica, agregamos por
\textbf{pooling estatístico} (média, desvio-padrão, percentil~90 e máximo) das
características de todos os cortes em um único vetor por exame --- decisão
justificada pela esparsidade espacial do achado, que descarta a média isolada e
motiva estatísticas de cauda capazes de capturar o corte mais anômalo. Alternativas
descartadas (características 2.5D e 3D diretas) são discutidas na Conclusão.

\subsection{Modelagem}
Implementamos o baseline trivial obrigatório (classificador de classe majoritária) e
seis famílias de modelos clássicos: regressão logística regularizada, SVM linear e
RBF \cite{cortes1995svm}, $k$-NN, Random Forest \cite{breiman2001randomforests},
\emph{gradient boosting} (XGBoost) \cite{chen2016xgboost} e um MLP raso (uma camada
oculta). Cada modelo é encapsulado em um \texttt{Pipeline} com padronização,
seleção/redução de dimensionalidade opcional e SMOTE opcional
\cite{chawla2002smote}, comparado contra o uso de \texttt{class\_weight='balanced'}.

\subsection{Protocolo experimental e métricas}
\label{sec:protocolo}
O particionamento é feito por \textbf{paciente} (nunca por corte ou por exame
isoladamente), via \texttt{StratifiedGroupKFold} com 5 dobras, semente fixa. O ajuste
de hiperparâmetros ocorre por busca em grade dentro de uma validação cruzada
aninhada (agrupada da mesma forma), nunca no conjunto de teste da dobra externa.
Reportamos AUC-ROC, AUC-PR, sensibilidade, especificidade, F1 e acurácia balanceada,
com média e desvio-padrão entre dobras --- acurácia isolada não é reportada como
métrica de desempenho, por ser enganosa em uma base com forte desbalanceamento de
classes \cite{colak2021rsnadataset}.

\section{Resultados}
\label{sec:resultados}

\noindent\textit{Nota metodológica (remover na versão final):} a
Tabela~\ref{tab:resultados} foi gerada por \texttt{scripts/run\_experiment.py} sobre
exames \textbf{sintéticos} (\texttt{src/data/synthetic.py}), pois este ambiente não
tinha acesso ao Kaggle. Ela comprova que o pipeline executa corretamente de ponta a
ponta, mas \textbf{não tem valor científico}. Regenerar sobre a amostra real de
\texttt{data/raw/}, conforme o \texttt{README.md}.

\begin{table}[h]
\centering
\caption{Comparação descritor~$\times$~modelo (média~$\pm$~desvio, 5 dobras). Exemplo sobre dados sintéticos de validação do pipeline --- substituir pelos resultados reais.}
\label{tab:resultados}
\small
\renewcommand{\arraystretch}{0.85}
\begin{tabular}{lccc}
\toprule
\textbf{Modelo} & \textbf{AUC-ROC} & \textbf{AUC-PR} & \textbf{Bal. Acc.} \\
\midrule
Baseline trivial   & 0,500 & 0,250 & 0,500 \\
Regressão logística & --     & --     & --     \\
SVM (linear/RBF)    & --     & --     & --     \\
$k$-NN              & --     & --     & --     \\
Random Forest        & --     & --     & --     \\
XGBoost              & --     & --     & --     \\
MLP raso             & --     & --     & --     \\
\bottomrule
\end{tabular}
\end{table}

\noindent As curvas ROC/PR e a matriz de confusão do melhor modelo (geradas
automaticamente em \texttt{outputs/figures/}) devem ser inseridas aqui a partir da
execução sobre dados reais, acompanhadas de análise de erro qualitativa (exemplos de
acerto e falha, com hipóteses de causa) e importância de características por
permutação.

\section{Conclusão}

Este trabalho estabeleceu a infraestrutura completa e reprodutível de um baseline
clássico para detecção de EP em angio-TC: leitura/pré-processamento DICOM
justificados, cinco famílias de descritores hand-crafted, sete modelos (incluindo o
baseline trivial) e um protocolo sem vazamento de dados. A limitação mais relevante é
a ausência, neste estágio, de segmentação vascular/pulmonar supervisionada -- usamos
recorte de margem fixa e limiar de intensidade como proxies de ROI, o que
provavelmente subestima o desempenho alcançável pelas características de forma e
radiômicas. Trabalho futuro imediato: (i) substituir os dados sintéticos pela amostra
real do RSNA STR PE Dataset e reportar os números definitivos na
Seção~\ref{sec:resultados}; (ii) avaliar características 2.5D como alternativa ao
pooling estatístico; (iii) comparar este piso clássico contra aprendizado profundo
(ex.: réplica simplificada do PENet \cite{huang2020penet}) na fase seguinte do
projeto.

\let\oldthebibliography\thebibliography
\renewcommand\thebibliography[1]{\oldthebibliography{#1}\setlength{\itemsep}{0pt}\setlength{\parskip}{0pt}}
{\footnotesize
\bibliographystyle{plain}
\bibliography{references}
}

\end{document}
